% Autoresearch Loops -- 15-minute defence deck
%   pdflatex presentation.tex            -> slides
%   uncomment \setbeameroption below     -> slides + speaker notes
\documentclass[aspectratio=169,11pt]{beamer}
\usetheme{metropolis}
\usepackage{appendixnumberbeamer}
\usepackage{booktabs}
\usepackage{pgfplots}
\usepgfplotslibrary{groupplots}
\pgfplotsset{compat=1.18}
\usetikzlibrary{shapes.geometric,arrows,positioning}

% \setbeameroption{show notes on second screen=right}   % presenter view
% \setbeameroption{show only notes}                     % notes handout

\metroset{block=fill,numbering=fraction,progressbar=frametitle}
\setbeamertemplate{note page}{\insertnote}

\title{Autoresearch Loops}
\subtitle{A Causal Study of the Instruction File}
\author{Christian Block}
\institute{Research project with Prof.\ Ramesh}
\date{\today}

\begin{document}
\maketitle

% ---------------------------------------------------------------- MOTIVATION
\begin{frame}{Nobody knows why these loops work}
\begin{itemize}
  \item March 2026: \texttt{autoresearch} ships
  \item ${\sim}13{,}000$ forks in six months
  \item An agent runs its own experiments
  \item Nobody can say \emph{which instruction} worked
\end{itemize}
\vfill
\begin{center}\small
\alert{Visual:} full-bleed screenshot of \texttt{github.com/karpathy/autoresearch},
star and fork counts visible
\end{center}
\note{In March 2026, Andrej Karpathy released a repository called autoresearch.
The idea is simple and slightly unsettling. You give a language model a goal,
some code it is allowed to edit, and a script that runs a trial and hands back
one number. Then you walk away. Overnight the model proposes a change, runs it,
reads the result, decides what to try next, and repeats until the budget is
gone. In six months it has been forked over thirteen thousand times. People are
running it on real systems. And here is the part that motivated this project:
when one of these loops produces a good result, nobody can tell you which part
of the instructions was responsible. That is the gap I set out to close.}
\end{frame}

\begin{frame}{The architecture is three files}
\begin{center}
\begin{tikzpicture}[font=\small,
  f/.style={rectangle,draw,thick,rounded corners=3pt,text width=3.1cm,
            align=center,minimum height=1.5cm,inner sep=4pt}]
\node[f,fill=gray!8]        (p) at (0,0)   {\texttt{prepare.py}\\[2pt]\scriptsize data + score\\\scriptsize\alert{nobody edits}};
\node[f,fill=gray!8]        (t) at (4.2,0) {\texttt{train.py}\\[2pt]\scriptsize the model\\\scriptsize\alert{agent edits}};
\node[f,fill=mLightBrown!25](m) at (8.4,0) {\texttt{program.md}\\[2pt]\scriptsize instructions\\\scriptsize\alert{human edits}};
\end{tikzpicture}
\end{center}
\vfill
\begin{itemize}
  \item No scheduler. No driver program.
  \item The cycle exists only as English prose
  \item \alert{\texttt{program.md} is not configuration --- it is the controller}
\end{itemize}
\note{The architecture is a contract between three files, and what matters is
who owns each one. prepare dot py fixes the data and defines the score. Nobody
touches it, so every experiment is scored identically. train dot py holds the
model, and the agent is the only thing that edits it. program dot md holds the
instructions, and the human is the only one who edits that. Now notice what is
missing. There is no scheduler. There is no driver program. The nine-step
experiment cycle exists only as English prose inside program dot md, and the
model itself is the execution layer. So program dot md is not configuration
handed to a controller. It IS the controller. That is the observation the whole
project rests on.}
\end{frame}

\begin{frame}{It already works in production. That is the problem.}
\begin{columns}[T]
\begin{column}{0.52\textwidth}
\begin{itemize}
  \item Shopify Liquid, PR \#2056
  \item \alert{53\%} faster parse and render
  \item ${\sim}120$ experiments, 93 commits
  \item Author: ``probably overfitted''
  \item \alert{Never merged}
\end{itemize}
\end{column}
\begin{column}{0.44\textwidth}
\begin{center}\small
\alert{Visual:} screenshot of the PR header\\with the \textsc{open} badge
\end{center}
\end{column}
\end{columns}
\note{This is not a toy concern. Here is an autoresearch run against Shopify's
Liquid template engine, which is production infrastructure. It cut parse and
render time by fifty-three per cent, over roughly a hundred and twenty
experiments across ninety-three commits. And then it was never merged. The
author's own assessment was that the gain was probably overfitted. So we have a
system that produced a large, real, measurable improvement, and the people
closest to it could not decide whether to trust it. Benchmarks have the same
shape. The AI Scientist, Coscientist, MLAgentBench --- all of them report
whether the pipeline succeeded end to end. None of them tell you which part of
the instruction file did the work. Success is measured; attribution is not.}
\end{frame}

\begin{frame}{Two questions}
\vfill
\begin{block}{RQ1}
Which components of the instruction file \alert{causally determine} which
aspects of the loop's behaviour, and how large is each effect?
\end{block}
\vfill
\begin{block}{RQ2}
Do the components act \alert{independently} of one another?
\end{block}
\vfill
\note{The obstacle is confounding. An instruction file settles several things at
once: what counts as success, how widely the agent may search, when it must
stop, how it must report, where to spend effort. If you compare two differently
worded files, every sentence differs, so nothing is identified. And the
prompt-sensitivity literature makes this sharper --- Sclar and colleagues showed
that formatting changes which preserve meaning can move accuracy by seventy-six
points. Wording clearly matters, and reading the wording will not tell you how.
So I ask two questions. First: which components causally determine which aspects
of the loop's behaviour, and how large is each effect? Second: do those
components act independently, or does the effect of one depend on another?}
\end{frame}

% --------------------------------------------------------------- METHODOLOGY
\begin{frame}{The loop under study}
\begin{center}
\begin{tikzpicture}[font=\scriptsize,node distance=1.5cm,
  b/.style={rectangle,draw,thick,rounded corners=2pt,minimum height=0.85cm,
            text width=2.05cm,align=center},
  llm/.style={b,fill=mLightBrown!30},
  det/.style={b,fill=gray!12},
  a/.style={->,thick}]
\node[llm] (p)                    {Propose};
\node[det] (e) [right=of p]       {Execute\\\scriptsize deterministic};
\node[llm] (i) [right=of e]       {Interpret};
\node[b]   (d) [right=of i]       {Decide};
\draw[a] (p)--(e); \draw[a] (e)--(i); \draw[a] (i)--(d);
\draw[a] (d.south) -- ++(0,-0.7) -| node[pos=0.25,below]{continue} (p.south);
\end{tikzpicture}
\end{center}
\begin{itemize}
  \item Shaded = model call \quad Unshaded = harness
  \item Search space: hyperparameters, not code
  \item \alert{The model never computes a number}
\end{itemize}
\note{This is the loop. The model alternates between two moves, each a single
JSON object. A propose carries a hypothesis, a model family, and hyperparameter
values. An interpret carries a reading of the last measurement and a decision to
continue or stop. If you have seen ReAct or the Ralph pattern, this is that
shape --- reason, act, observe, repeat. One narrowing matters. The search space
is a hyperparameter space, not arbitrary code, so that the instruction file
stays the only free variable. And critically, the model never computes a number.
It picks what gets computed, and reads the answer back. The shaded boxes are
model calls; the unshaded one is a deterministic harness running scikit-learn.
That means every value in a transcript is traceable: either the executor
produced it, and it is a fact about the task, or the model produced it, and it
is a fact about the loop's behaviour.}
\end{frame}

\begin{frame}{Stop reading the prompt. Start configuring it.}
\begin{columns}[T]
\begin{column}{0.5\textwidth}
\begin{itemize}
  \item One template, five variable slots
  \item Everything else byte-identical
  \item Each slot written at two levels
  \item \alert{A file becomes a vector}
\end{itemize}
\end{column}
\begin{column}{0.46\textwidth}
\begin{center}
\vspace{0.3cm}
$c=(c_1,\ldots,c_5),\quad c_i\in\{0,1\}$\\[6pt]
{\large$c=(1,0,1,0,1)$}\\[4pt]
{\scriptsize $2^5=32$ distinct instruction files}
\end{center}
\end{column}
\end{columns}
\note{Here is the core move of the project. Instead of treating the instruction
file as prose to be interpreted, I treat it as a configuration vector. Every
file comes out of one template. The fixed part --- the research question, the
executor interface, the response format --- is byte-identical across every file
I generate. The variable part is exactly five paragraphs, each written at one of
two levels. Nothing else in the file moves. So a complete instruction file is
described by five binary digits. That is what makes the problem tractable: once
a prompt is a vector, the question what does this file do becomes a question you
can answer with an experiment instead of an argument.}
\end{frame}

\begin{frame}{The five components}
\small
\begin{tabular}{@{}clll@{}}
\toprule
& \textbf{Component} & \textbf{Level 0} & \textbf{Level 1}\\
\midrule
$M$ & Evaluation criterion & undefined      & 5-fold CV accuracy\\
$B$ & Search breadth       & one family     & both families\\
$S$ & Stopping rule        & fixed 3        & adaptive, up to 8\\
$O$ & Reporting style      & terse          & full sentences\\
$E$ & Search emphasis      & explore first  & exploit first\\
\bottomrule
\end{tabular}
\vfill
\begin{itemize}\small
  \item Categories taken from a real \texttt{program.md}
  \item All ten paragraphs LLM-written, then frozen
\end{itemize}
\note{These are the five components. Evaluation criterion: either success is
left undefined, or it is pinned to five-fold cross-validated accuracy. Search
breadth: commit to one model family, or compare both. Stopping rule: a fixed
budget of exactly three experiments, or an adaptive budget of up to eight that
halts on diminishing returns. Reporting style: terse, or full sentences stating
hypothesis and interpretation. And emphasis: explore several configurations
first, or exploit a promising one immediately. These are not invented for the
study. A real program dot md hardcodes a baseline, says how to handle failures,
fixes a reporting discipline, and imposes a parsimony rule. These five come from
those categories. One control worth noting: a language model wrote all ten
paragraphs from fixed meta-prompts, and they were then frozen and reused word
for word --- so my own stylistic habits are not a free parameter tangled up in
the treatment.}
\end{frame}

\begin{frame}{Run the whole space}
\begin{center}
{\Large $N \;=\; \underbrace{2^5}_{\text{32 configs}} \times \underbrace{D}_{\text{2 tasks}} \times \underbrace{R}_{\text{20 reps}} \;=\; 1{,}280$ runs}
\end{center}
\vfill
\begin{itemize}
  \item Constructed, not observed
  \item Confounding gone \alert{by design}, not adjustment
  \item Every interaction estimable from every main effect
  \item A fraction would have to \alert{assume} RQ2's answer
\end{itemize}
\note{I run the complete space. All thirty-two configurations, crossed with two
datasets, twenty replicates each --- twelve hundred and eighty real runs against
a real executor. Two things follow from running all of it. First, because I
construct the configurations rather than collecting them, confounding is
eliminated by the design itself, with no statistical adjustment. In
observational prompt data, whoever writes a precise success criterion probably
also writes a careful stopping rule --- those two would be correlated.
Construction breaks that. Second, with all thirty-two files in hand, every
pairwise interaction is estimated separately from every main effect. A
fractional design would have to assume the components act independently, because
the runs that distinguish a main effect from an interaction are exactly the ones
it skips. Thirty-two cells is the price of not assuming, and at this scale I can
pay it.}
\end{frame}

\begin{frame}{What I held fixed --- and one thing I did not}
\begin{columns}[T]
\begin{column}{0.48\textwidth}
\textbf{Held constant}
\begin{itemize}\small
  \item Pinned model snapshot
  \item One stratified 70/30 split
  \item Same executor and whitelist
  \item Same seed list in every cell
\end{itemize}
\end{column}
\begin{column}{0.48\textwidth}
\begin{alertblock}{Temperature $=0.7$, not $0$}
\small Run-to-run variance \emph{is} the error term.\\
Greedy decoding would zero it --- and take
every standard error with it.
\end{alertblock}
\end{column}
\end{columns}
\note{Controls. The agent model is pinned to a dated snapshot rather than a
moving alias. The split, the seed, the executor, the whitelist, the fixed part
of the template --- all constant. Runs are separate conversations, so nothing
carries over. Now the one that gets questioned, so let me address it directly. I
did not run at temperature zero. It is fixed at zero point seven for every
single run. This is deliberate. Run-to-run variance is the error term that every
coefficient in this study is tested against. A deterministic loop would drive
that variance to zero and take every standard error with it --- I would have
point estimates and no way to say whether any of them were distinguishable from
noise. So I hold temperature constant and above zero, which makes it a
controlled constant rather than an eliminated one, and I pass an explicit
sampling seed so any individual run can still be regenerated exactly.}
\end{frame}

\begin{frame}{Four metrics, straight off the transcript}
\begin{itemize}
  \item \textbf{Gain over default} --- did it improve anything?
  \item \textbf{Wasted trial ratio} --- malformed $+$ duplicate
  \item \textbf{Improvement rate} --- trials beating the best so far
  \item \textbf{Cost to best} --- calls to reach the peak
\end{itemize}
\vfill
\begin{center}
\alert{No judge model. No scoring library. No evaluation framework.}
\end{center}
\note{Four dependent variables, all computed mechanically from the structured
transcript. Gain over default: the score of the configuration the run finally
recommends, minus a default-hyperparameter baseline computed before any run.
Wasted trial ratio: proposals that were malformed or rejected by the whitelist,
plus exact duplicates of earlier proposals, over all attempts --- both of those
burn a turn without buying a measurement. Improvement rate: the share of
executed trials that beat the best score seen earlier in the same run. And cost
to best: how many executor calls it took to reach the peak. No evaluation
framework, no scoring library, no judge model is involved anywhere. Every metric
is a difference in accuracy, a share of proposals, or a count of calls.}
\end{frame}

\begin{frame}{The model --- and why the dataset is in it}
\begin{center}
{\large $Y \;=\; \beta_0 \;+\; \sum_{i=1}^{5}\beta_i X_i \;+\; \alert{\sum_{j=2}^{D}\gamma_j D_j} \;+\; \varepsilon$}
\end{center}
\vfill
\begin{itemize}
  \item $\pm1$ coding --- orthogonal by construction
  \item $\gamma_j$ absorbs \alert{task difficulty}
  \item Without it, that variance inflates every standard error
  \item HC3 errors; Benjamini--Hochberg for multiplicity
\end{itemize}
\note{Each component is coded plus or minus one, and I regress each metric on
those five indicators --- plus a dataset fixed effect, the gamma term
highlighted here. That term is not decoration. Breast Cancer and Wine differ in
how much headroom they leave and how they respond to the whitelist. Without the
dataset in the model, that variation sits in the residual, inflates every
standard error, and makes real effects harder to see. Because the factorial is
crossed with the task suite, each component indicator is orthogonal to the
dataset indicator by construction, so the coefficients I report are within-task
effects. I fit on individual runs rather than cell means, use
heteroskedasticity-consistent standard errors --- the stopping rule pins the
experiment count at one level and frees it at the other, so the variance
genuinely is not constant --- and apply Benjamini-Hochberg for multiplicity.}
\end{frame}

% ------------------------------------------------------------------- RESULTS
\begin{frame}{What actually moved}
\begin{center}\begin{tikzpicture}
\begin{axis}[width=10.5cm,height=5.4cm,y dir=reverse,
  ytick={1,2,3,4,5}, yticklabels={Evaluation criterion,Search breadth,Stopping rule,Reporting style,Search emphasis},
  ymin=0.5,ymax=5.5,xmin=-0.1001,xmax=0.1001,
  xlabel={change in wasted trial ratio ($2\hat\beta$)},
  ymajorgrids,major grid style={gray!25,dotted},
  axis line style={gray!60},tick align=outside,tick pos=left,
  xtick={-0.05,0,0.05},scaled x ticks=false,
  xticklabel style={/pgf/number format/fixed,/pgf/number format/precision=2},]
\addplot[gray!70,dashed,forget plot] coordinates {(0,0.5) (0,5.5)};
\addplot[only marks,mark=*,mark size=3pt,mLightBrown,very thick,
  error bars/.cd,x dir=both,x explicit,error bar style={mLightBrown,thick}]
  coordinates {(-0.02318,1) +- (0.01218,0) (-0.03370,2) +- (0.01218,0) (0.04396,3) +- (0.01218,0) (-0.07484,4) +- (0.01218,0)};
\addplot[only marks,mark=o,mark size=3pt,gray,very thick,
  error bars/.cd,x dir=both,x explicit,error bar style={gray,thick}]
  coordinates {(0.00490,5) +- (0.01218,0)};
\end{axis}
\end{tikzpicture}
\end{center}
\vspace{-0.2cm}
\begin{center}\scriptsize
filled $=$ survives Benjamini--Hochberg \quad bars $=$ 95\% CI
\end{center}
\note{Here is what actually moved. The wasted trial ratio is the metric the
instruction file really governs --- four of the five components shift it, and
the model explains sixteen per cent of its variance. Reporting style produces
the largest effect anywhere in the design: minus zero point zero seven five. In
plain terms, requiring the model to write full sentences cuts wasted trials from
about ten and a half per cent to three per cent. That is a two-thirds reduction,
from a change in wording alone. The stopping rule is the only component that
touches either outcome metric --- worth about one tenth of a point on the gain,
and it lifts the improvement rate from thirteen and a half to seventeen per
cent. And the cost to best answers to nothing: no component survives adjustment.
The largest candidate, exploit-first, lands at p equals zero point zero five
seven. So close, and I am not going to claim it.}
\end{frame}

\begin{frame}{The task matters more than the instructions}
\begin{center}\begin{tikzpicture}
\begin{axis}[ybar,width=10cm,height=5.2cm,
  symbolic x coords={Improvement rate,Cost to best},
  xtick=data,ymin=0,bar width=22pt,enlarge x limits=0.5,
  ylabel={$|$effect$|$},legend style={at={(0.5,-0.22)},anchor=north,legend columns=2,draw=none},
  ymajorgrids,major grid style={gray!25,dotted},axis line style={gray!60},
  scaled y ticks=false,ytick={0,0.05,0.1,0.15,0.2},
  yticklabel style={/pgf/number format/fixed,/pgf/number format/precision=2},
  nodes near coords,nodes near coords style={font=\scriptsize},
  every node near coord/.append style={/pgf/number format/precision=3,/pgf/number format/fixed},]
\addplot[fill=gray!45,draw=gray!60] coordinates {(Improvement rate,0.0341) (Cost to best,0.1094)};
\addplot[fill=mLightBrown,draw=mLightBrown] coordinates {(Improvement rate,0.0638) (Cost to best,0.2031)};
\legend{largest component effect, dataset effect}
\end{axis}
\end{tikzpicture}
\end{center}
\vspace{-0.1cm}
\begin{itemize}\small
  \item Both dataset effects at $p<10^{-5}$
  \item \alert{Adaptive stopping rule obeyed in only 20\% of runs}
\end{itemize}
\note{Two findings worth dwelling on. First: on two of the four metrics, the
dataset effect is roughly twice the largest instruction component. Moving from
Breast Cancer to Wine shifts the improvement rate by six and a half points and
the cost to best by two tenths of a call, both at p below ten to the minus five.
Which task you pick matters more than how you word the prompt. That is exactly
the variation the fixed effect removes from the residual. Second, and this
constrains everything else: I measured whether the model actually obeyed each
instruction. The fixed budget is followed in eighty-nine per cent of runs. The
adaptive rule, in twenty. So the stopping-rule effects I just showed are the
effect of offering an adaptive budget, not of a loop that reliably stops on
diminishing returns. Every estimate here is intention-to-treat, attenuated
toward zero.}
\end{frame}

\begin{frame}{What we proved}
\vfill
\begin{itemize}
  \item Instructions govern \alert{process}, not outcome
  \item Verbose reporting: wasted trials $10.6\% \to 3.1\%$
  \item Components act independently (3 of 40 interactions)
  \item The method transfers to any prompt you can slice
\end{itemize}
\vfill
\begin{center}\small
$a_0 = 0.960$ on both tasks --- the ceiling, not the loop, is the binding constraint
\end{center}
\note{So, what did we prove. The instruction file governs how efficiently the
loop conducts its search, not the quality of what it finds. Requiring the model
to explain itself removes two thirds of the wasted proposals --- a large and
reproducible effect. But no wording tested here made the loop find a better
model, and I would argue that is a ceiling effect: untuned defaults already
reach ninety-six per cent. On RQ2, three of forty interactions survive
adjustment, all on one metric and all involving reporting style, so the
components can be treated as independent --- and because the factorial is
complete, that is something I tested rather than assumed. Finally, the method is
not tied to the five components I happened to vary. Any instruction file that
can be cut into slots can be studied this way. Thank you. I am happy to take
questions.}
\end{frame}

\appendix
\begin{frame}[standout]Questions\end{frame}

\begin{frame}{Backup --- anticipated questions}
\footnotesize
\begin{tabular}{@{}p{4.6cm}p{7.4cm}@{}}
\toprule
Why not temperature 0? & It would zero the error term every coefficient is tested against.\\
$R^2 = 0.16$ --- weak? & Expected for 5 factors over noisy agent behaviour. Hence $N=1280$.\\
Only one agent model? & Stated limitation: one policy, not instructions in general.\\
One wording per level? & Stated limitation: effect confounded with phrasing. Top follow-up.\\
Why no compliance for $M$? & Executor scores regardless, so $M$ leaves no trace. $O$ gets 205 vs 386 chars.\\
Only 40\% beat baseline? & Ceiling effect: $a_0=0.960$ on both tasks.\\
\bottomrule
\end{tabular}
\end{frame}

\begin{frame}{Backup --- main effects}
\centering\footnotesize
\begin{tabular}{lrrrrrrr}
\toprule
Metric & $M$ & $B$ & $S$ & $O$ & $E$ & $\hat\gamma$ & $R^2$\\
\midrule
Gain over default   & +0.001 & -0.001 & +0.001$^{**}$ & +0.000 & +0.000 & -0.001$^{***}$ & 0.029 \\
Wasted trial ratio  & -0.023$^{***}$ & -0.034$^{***}$ & +0.044$^{***}$ & -0.075$^{***}$ & +0.005 & -0.003 & 0.159 \\
Improvement rate    & +0.008 & -0.026$^{*}$ & +0.034$^{**}$ & +0.006 & +0.018 & -0.064$^{***}$ & 0.044 \\
Cost to best        & +0.06 & -0.02 & +0.05 & -0.01 & -0.11 & -0.20$^{***}$ & 0.025 \\
\bottomrule
\end{tabular}
\end{frame}

\end{document}
